\documentclass[10pt]{article}
%% korr-preamble.tex -- gemeinsame Präambel für die Korrekturlisten LNM 1184
%% Makros wie im Buch (LN-Book-Springer/preamble), damit die Formeln
%% genauso aussehen wie im Proof.
\usepackage[T1]{fontenc}
\usepackage[utf8]{inputenc}
\usepackage{newtxtext}
\usepackage[varvw]{newtxmath}
\usepackage[scaled=.9]{helvet}
\usepackage[varqu,scaled=.95]{zi4}
\usepackage[ngerman]{babel}
\usepackage{csquotes}
\usepackage[a4paper,landscape,margin=1.6cm,top=1.8cm,bottom=1.8cm]{geometry}
\usepackage{microtype}
\usepackage{array,longtable,booktabs}
\usepackage{ragged2e}
\usepackage[table]{xcolor}
\usepackage{xspace}
\usepackage{enumitem}
\setlist{nosep,leftmargin=*}
\usepackage[nointegrals]{wasysym}
\usepackage{newunicodechar}
\usepackage{nicefrac}
\usepackage{mathtools}
\usepackage{mleftright}
\mleftright
\usepackage[hidelinks]{hyperref}
\setlength{\parindent}{0pt}
\setlength{\parskip}{.5em}
\setlength{\emergencystretch}{3em}
\renewcommand{\arraystretch}{1.25}
\definecolor{codegray}{gray}{0.35}

%% ---- Makros aus dem Buch -------------------------------------------
\DeclarePairedDelimiterX{\norm}[1]{\lVert}{\rVert}{\ifblank{#1}{\:\cdot\:}{#1}}
\DeclarePairedDelimiterX{\abs}[1]{\lvert}{\rvert}{\ifblank{#1}{\:\cdot\:}{#1}}
\DeclarePairedDelimiterX{\dualp}[1]{\langle}{\rangle}{\ifblank{#1}{\,\cdot\,,\cdot\,}{#1}}
\newcommand{\interval}[1]{\left[ #1 \right]}
\newcommand{\ointerval}[1]{\left] #1 \right[}
\newcommand{\rointerval}[1]{\left[ #1 \right[}
\newcommand{\lointerval}[1]{\left] #1 \right]}
\newcommand{\tint}{\textstyle\int}
\AtBeginDocument{%
\let\ORGvarepsilon=\varepsilon
\let\epsilon=\epsilon
\let\epsilon=\ORGvarepsilon
\let\ORGvarrho=\varrho
\let\rho=\rho
\let\rho=\ORGvarrho
\let\ORGvartheta=\vartheta
\let\theta=\theta
\let\theta=\ORGvartheta
\let\ORGvarphi=\varphi
\let\phi=\phi
\let\phi=\ORGvarphi
\let\ORGvarleq=\leqslant
\let\leqslant=\leq
\let\leq=\ORGvarleq
\let\ORGvargeq=\geqslant
\let\geqslant=\geq
\let\geq=\ORGvargeq
\let\ORGvarpreccurlyeq=\preccurlyeq
\let\preccurlyeq=\preceq
\let\preceq=\ORGvarpreccurlyeq
\let\setminus\smallsetminus			% A - B fuer Mengen
\renewcommand{\P}{\mathfrak{P}}		% Achtung: \P hat andere Bedeutung in LaTeX
}
\newcommand{\N}{\mathbb{N}}		% Natuerliche Zahlen
\newcommand{\Z}{\mathbb{Z}}		% Ganze Zahlen
\newcommand{\Q}{\mathbb{Q}}		% Rationale Zahlen
\newcommand{\R}{\mathbb{R}}		% Reelle Zahlen
\newcommand{\C}{\mathbb{C}}		% Komplexe Zahlen
\newcommand{\K}{\mathbb{K}}		% Koerperzeichen
\newcommand{\T}{\mathbb{T}}		% Torus
\DeclareMathOperator{\Id}{Id}
\DeclareMathOperator{\sign}{sign}				% signum-Operator
\newcommand{\im}{\ensuremath{\mathrm{i}}}		% imaginäre Konstante i
\newcommand{\iu}{\ensuremath{\mathrm{i}}}
\newcommand{\eu}{\ensuremath{\mathrm{e}}}
\renewcommand*{\Re}{\mathop{}\!\mathrm{Re}}		% Realteil
\renewcommand*{\Im}{\mathop{}\!\mathrm{Im}}		% Imaginärteil
\newcommand*{\supp}{\mathop{}\!\mathrm{supp}\,}  % Träger
\newcommand*{\diff}[1]{\mathop{}\!\mathrm{d}{#1}}	% \diff{s} = ds, 
\newcommand*{\ds}{\mathop{}\!\mathrm{d}{s}}       	% \ds = ds, 
\newcommand*{\dg}{\mathop{}\!\mathrm{d}{g}}       	% \dg = dg, 
\newcommand*{\dt}{\mathop{}\!\mathrm{d}{t}}       	% \dt = dt, 
\newcommand*{\dx}{\mathop{}\!\mathrm{d}{x}}       	% \dx = dx, 
\newcommand*{\dy}{\mathop{}\!\mathrm{d}{y}}       	% \dy = dy, 
\newcommand*{\dr}{\mathop{}\!\mathrm{d}{r}}       	% \dr = dr, 
\newcommand*{\dm}{\mathop{}\!\mathrm{d}{m}}       	% \dm = dm, 
\newcommand*{\du}{\mathop{}\!\mathrm{d}{u}}       	% \du = du, 
\DeclareMathOperator{\esssup}{ess\,sup}
\newcommand{\LE}{\mathcal{L}(E)} 		% 
\newcommand{\BH}{\mathcal{B}(H)} 		% 
\newcommand{\BA}{\mathfrak{A}}
\newcommand{\ZE}{\mathcal{Z}(E)} 	% \ZE  Zentrum von E
\newcommand{\LH}{\mathcal{L}(H)} 
\let\originalL\L
\AtBeginDocument{%
  \DeclareRobustCommand{\L}[1]{\mathcal{L}(#1)}%
}
\newcommand{\Fix}[1]{\mathop{}\mathrm{Fix}\left(#1\right)}		% \Fix{T} Fixraum
\newcommand{\rank}{\mathop{}\!\mathrm{rank}\,}		% 
\newcommand{\Kern}[1]{\mathop{}\mathrm{ker}(#1)}
\newcommand{\Image}[1]{\mathop{}\mathrm{im}(#1)}
\newcommand{\image}{\mathop{}\!\mathrm{im}\,}
\renewcommand{\dim}{\mathop{}\!\mathrm{dim}\,}
\newcommand*{\eg}{e.g.,\xspace}
\newcommand*{\ie}{i.e.,\xspace}
\newcommand*{\resp}{resp.\xspace}
\newcommand*{\vs}{vs.\xspace}
\newcommand*{\cf}{cf.\xspace}
\newcommand*{\etc}{etc.\xspace}
\newcommand{\CA}{$\mathrm{C}^{*}$}	% C*-Algebra
\newcommand{\WA}{$\mathrm{W}^{*}$}	% W*-Algbera
\newcommand{\AW}{$\mathrm{AW}^{*}$}	% AW*-Algbera
\newcommand*{\TT}{\mathcal{T}}		% C_0-Semigroup
\newcommand*{\RR}{\mathcal{R}}		% Pseudo Resolvente
\newcommand*{\SG}{\mathcal{S}}	    % Semigroup
\newcommand*{\UG}{\mathcal{U}}	    % Implemented Semigroup
\newcommand*{\F}{\mathcal{F}}		% F-product
\newcommand{\mfrac}[2]{%
    {\raise.5ex\hbox{$\scriptstyle #1$}}%
    \kern.025em/\kern.025em%
    {\lower.25ex\hbox{$\scriptstyle #2$}}}
\newcommand*{\sfrac}[2]{\leavevmode\kern.1em%
       \raise.5ex\hbox{\small #1}%
       \kern.025em/\kern.025em%
       \lower.25ex\hbox{\small #2}}
\newcommand{\1}{\mathbf{1}}
\newcommand{\HO}{\mathcal{H}(O)} 		%  für Lotz
\newcommand{\WO}{\mathcal{W}(O)} 		%  für Lotz
\newcommand*{\ddp}{\mathop{}\!\mathrm{d}{p}}       	% \ddp = dp, 
\newcommand*{\dN}{\mathop{}\!\mathrm{d}{N}}       	% \dN = dN, 
\newcommand*{\ddt}{\mathop{}\!\mathrm\fraq{d}{dt}}   % \ddt = d/dt
\newcommand*{\lnm}{LNM1184: } % \LMN-BookCode
\newcommand{\closure}[1]{#1^{-}}
%% Hilfsmatrizen aus Teil I, Kap. 1
\newcommand{\ABC}{\begin{pmatrix} 0 & 1 \\ -1 & 0 \end{pmatrix}}
\newcommand{\BBC}{\begin{pmatrix} \cos(t) & \sin(t) \\ - \sin(t) & \cos(t) \end{pmatrix}}
\newcommand{\CBC}{\begin{pmatrix} 1 & 1 \\ -1 & -1 \end{pmatrix}}
\newcommand{\DBC}{\begin{pmatrix} 1+t & t \\ - t & 1-t \end{pmatrix}}
\newcommand{\EBC}{\begin{pmatrix} 0 & 1 \\ 1 &0 \end{pmatrix}}
\newcommand{\FBC}{\begin{pmatrix} \cosh(t) & \sinh(t) \\ \sinh(t) & \cosh(t) \end{pmatrix}}

%% ---- Darstellung der Listen ----------------------------------------
%% Quelltext, der sich nicht als Formel setzen lässt (Labels, \ref, Dateinamen)
\newcommand{\code}[1]{\leavevmode{\color{codegray}\ttfamily\small #1}}
\newcommand{\sichA}{\CIRCLE}      % ● hoch
\newcommand{\sichB}{\LEFTcircle}  % ◐ mittel
\newcommand{\sichC}{\Circle}      % ○ niedrig
%% Querverweise im Buch: Label-Namen statt Nummern
\newcommand{\lbl}[1]{\textsf{\footnotesize\color{codegray}\detokenize{#1}}}

\hypersetup{pdftitle={Formelprüfung Part C (Teil III): Satz und Notation}}
\pagestyle{plain}
\begin{document}
% Erzeugt aus korrekturen-mathe-part-c.md: Formeln gesetzt statt Quelltext.
% Grau/Schreibmaschine = Quelltext, der sich nicht als Formel setzen laesst (Makronamen, Labels).
{\Large\bfseries Formelprüfung Part C (Teil III): Satz und Notation\par}
\smallskip
\emph{One-parameter Semigroups of Positive Operators}, 2. Aufl. (Springer LNM 1184). Geprüft am 14.10.2026; Stand der Umsetzung: 20.10.2026.

\textbf{Umfang:} Teil III, Kap. 1–4, geprüft im LaTeX-Quelltext (Branch \code{fahnenkorrektur}). Geprüft wurden nur Satz und Notation in Formeln: Tippfehler, Klammern, Indizes, uneinheitliche Schreibweisen, Mathemodus, Satzzeichen an Formeln, Verweise. Beweise wurden \textbf{nicht} nachgerechnet; wo ein Befund inhaltlich wirkt, steht „bitte prüfen“.

\textbf{Stand 20.10.2026:} \textbf{Kap. 1 und Kap. 2 sind erledigt} (Kap. 1 nach Durchsicht durch Rainer Nagel; in der Spalte Kommentar mit \textbf{Umgesetzt} markiert, Abweichungen ebenda). Kap. 3–4 sind noch offen. Die Zeilenangaben beziehen sich auf den Stand vom 14.10.2026.

\textbf{Buchführung:} Maßgeblich für den Korrekturstand sind nur die Listen {korrekturen-text-*}, {korrekturen-mathe-*} und {math-error-*} (seit 20.10.2026). \code{Umgesetzte\_Korrekturen\_LNM1184.md}, \code{Finales-zum-Pruefen.md} und \code{Fahnenkorrektur\_Springer\_LNM1184.md} werden nicht mehr fortgeführt.

\textbf{Spalten:} S. $=$ gedruckte Buchseite, ermittelt über das nächstgelegene Label (kann um eine Seite abweichen). Z. $=$ Zeile in der .tex-Datei. Sicherheit: \sichA{} eindeutig $\ensuremath{\cdot}$ \sichB{} sehr wahrscheinlich $\ensuremath{\cdot}$ \sichC{} Hinweis/Ermessen.

Insgesamt \textbf{222 Punkte}: 47 \sichA{}, 99 \sichB{}, 76 \sichC{}. Davon stehen 26 bereits in der Arbeitsliste \code{Fahnenkorrektur\_Springer\_LNM1184.md} (im Kommentar markiert).

\section*{Einzelbefunde}
\subsection*{Teil III, Kap. 1 – Basic Results on Banach Lattices and Positive Operators (\code{part-c/chap-c1.tex})}
\begin{longtable}{@{}rrc>{\RaggedRight\arraybackslash}p{5.8cm}>{\RaggedRight\arraybackslash}p{5.8cm}>{\RaggedRight\arraybackslash}p{1.9cm}>{\RaggedRight\arraybackslash}p{7.3cm}@{}}
\toprule
\textbf{S.} & \textbf{Z.} & \textbf{} & \textbf{Ist} & \textbf{Soll} & \textbf{Art} & \textbf{Kommentar} \\
\midrule
\endfirsthead
\toprule
\textbf{S.} & \textbf{Z.} & \textbf{} & \textbf{Ist} & \textbf{Soll} & \textbf{Art} & \textbf{Kommentar} \\
\midrule
\endhead
\bottomrule
\endlastfoot
265 & 35 & \sichB{} & $\mathrm{(LO1)} \quad$ & $\text{(LO$_{1}$)} \quad$ & Notation & \textbf{Umgesetzt.} In der Definition „(LO1)“/„(LO2)“ ohne Index, im Fließtext Z. 45, 46, 49 „(LO$_{1}$)“/„(LO$_{2}$)“ mit tiefgestellter Ziffer, im Druck sichtbar verschieden; ebenso Z. 37 $\mathrm{(LO2)}$. Alternativ Z. 45–49 auf (LO1)/(LO2) umstellen. \\[2pt]
265 & 71 & \sichA{} & $\inf   \sup_{\alpha} f_{\alpha},h  )$ & $\inf(\sup_{\alpha} f_{\alpha},h)$ & Klammern & \textbf{Umgesetzt.} Öffnende Klammer nach \code{\textbackslash{}inf} fehlt (vgl. die Parallelformel Z. 72 $\sup(\inf_{\alpha}f_{\alpha},h)$). \\[2pt]
266 & 79 & \sichC{} & \code{\textbackslash{}begin\{equation*\}\allowbreak{}\textbackslash{}label\{eq:c1-1\}} & \code{\textbackslash{}begin\{equation*\}} & LaTeX & \textbf{Umgesetzt.} \code{\textbackslash{}label} in einer unnummerierten Umgebung; das Label wird nirgends referenziert (ein \code{\textbackslash{}ref} würde die Abschnittsnummer liefern). Ebenso Z. 267 (eq:c1-2) und Z. 275 (eq:c1-3). Im Druck unsichtbar, nur Bereinigung. \\[2pt]
269 & 206 & \sichC{} & $\{u\}^{d} = 0$ & $\{u\}^{d} = \{0\}$ & Notation & \textbf{Umgesetzt.} Gleichheit von Mengen; sonst im Kapitel $= \{0\}$ (Z. 180, 368). \\[2pt]
270 & 245 & \sichB{} & $\sup\{\phi(h) \colon \vert h\vert \leq f \text{ for $ f \in E_{+} $}\}$ & $\sup\{\phi(h) \colon \vert h\vert \leq f\} \quad \text{for $ f \in E_{+} $}$ & Notation & \textbf{Umgesetzt.} Die Bedingung „for $f \ensuremath{\in} E\ensuremath{_{+}}$“ steht innerhalb der Mengenklammer, betrifft aber die feste Variable f, nicht die Laufvariable h; gehört hinter die Menge. Ebenso Z. 384 ($\vert h\vert \leq f , \text{ $ f \in E_{+} $} \}$) und Z. 529 ($\text{ for $ h \geq 0 $} \}$). \\[2pt]
271 & 276 & \sichC{} & $g \in E_{+}'$ & $g \in E'_{+}$ & Notation & \textbf{Umgesetzt.} Strich hinter dem Index gesetzt; in Teil III sonst durchweg $E'_{+}$ (z. B. chap-c2 Z. 61, 415, 488, 568). \\[2pt]
271 & 302 & \sichB{} & {which we denote by $ (E,x') $} & {which we denote by $ (E,f') $} & Tippfehler & \textbf{Umgesetzt.} Das Funktional heißt im Kontext f' $(q_{f'})$, ein x' kommt nicht vor; vgl. Z. 305 $(E,\mu)$. Bitte prüfen. \\[2pt]
273 & 371 & \sichB{} & $\mathcal{L}^{r}(E;F)$ & $\mathcal{L}^{r}(E,F)$ & Notation & \textbf{Umgesetzt.} Semikolon als Trenner, während im selben Absatz $\L{E,F}$ mit Komma steht (Z. 369, 372, 373, 475); sichtbar uneinheitlich. Ebenso Z. 372, 373, 374, 381. \\[2pt]
273 & 435 & \sichA{} & $h = f + ig$ & $h = f + \iu g$ & Notation & \textbf{Umgesetzt.} Kursives i für die imaginäre Einheit; sonst \code{\textbackslash{}iu/\allowbreak{}\textbackslash{}im} (aufrecht), z. B. Z. 428, 438, 446, 450, 461. Ebenso Z. 474 $T_{1} + iT_{2}$ und Z. 533 $+ i(S_{3} - S_{4})$. \\[2pt]
273 & 476 & \sichB{} & {if\/ $ T_{2} $ is $ = 0 $} & {if\/ $ T_{2} = 0 $} & Mathemodus & \textbf{Umgesetzt} (stand schon so im Text). Schon in der Arbeitsliste. Relationszeichen allein im Mathemodus nach „is“. Ähnlich Z. 524 {one has necessarily the values $ = 0 $ for $ S_{f} $} (besser: „one has necessarily $ S_{f} = 0 $“). \\[2pt]
275 & 508 & \sichC{} & $\text{sign } f = f\cdot\vert f\vert^{-1}$ & $\sign f = f\cdot\vert f\vert^{-1}$ & Notation & \textbf{Umgesetzt.} In Z. 485 mit dem Operator \code{\textbackslash{}sign} gesetzt; \code{\textbackslash{}text\{sign \allowbreak{}\}} übernimmt die Umgebungsschrift und hat anderen Abstand. \\[2pt]
275 & 520 & \sichB{} & {the band $ \{x\}^{dd} $} & {the band $ \{f\}^{dd} $} & Tippfehler & \textbf{Umgesetzt} (vorgezogen zusammen mit math-error Z. 529; Rainer ohne Einwand). Ein Element x kommt im Kontext nicht vor, gemeint ist das Band zu f. Ebenso Z. 524 $\{x\}^{d}$ und Z. 527 $\{x\}^{dd}$. Bitte prüfen. \\[2pt]
275 & 529 & \sichC{} & $S_{f}h = \sup \{ S_{f}g \colon$ & $S_{f}h = \sup \{ S_{0}g \colon$ & Tippfehler & \textbf{Umgesetzt} (stand schon so im Text). $S_f$ wird hier erst definiert; rechts ist wohl $S_0$ gemeint ($g \ensuremath{\in} E_{\vert f\vert}$). Bitte prüfen. \\[2pt]
276 & 579 & \sichA{} & {then $ Z(E) $ is isomorphic} & {then $ \mathcal{Z}(E) $ is isomorphic} & Notation & \textbf{Umgesetzt.} Zentrum sonst kalligraphisch $\mathcal{Z}(E)$ (Z. 572, 573, 575, 582). \\
\end{longtable}
\subsection*{Teil III, Kap. 2 – Characterization of Positive Semigroups on Banach Lattices (\code{part-c/chap-c2.tex})}
\begin{longtable}{@{}rrc>{\RaggedRight\arraybackslash}p{5.8cm}>{\RaggedRight\arraybackslash}p{5.8cm}>{\RaggedRight\arraybackslash}p{1.9cm}>{\RaggedRight\arraybackslash}p{7.3cm}@{}}
\toprule
\textbf{S.} & \textbf{Z.} & \textbf{} & \textbf{Ist} & \textbf{Soll} & \textbf{Art} & \textbf{Kommentar} \\
\midrule
\endfirsthead
\toprule
\textbf{S.} & \textbf{Z.} & \textbf{} & \textbf{Ist} & \textbf{Soll} & \textbf{Art} & \textbf{Kommentar} \\
\midrule
\endhead
\bottomrule
\endlastfoot
280 & 57 & \sichC{} & {(see Eq.~(\lbl{eq:a2-2.12}))} & {(see Part-I, Eq.~(\lbl{eq:a2-2.12}))} & Verweis & \textbf{Umgesetzt.} Formel aus Teil I ohne „Part-I,“ zitiert; das Kapitel hat eigene Formeln (2.2.x), die Nummer ist nur über die Seitenangabe eindeutig. Ebenso Z. 64 (eq:a2-2.5). \\[2pt]
280 & 61 & \sichB{} & $\partial N^{+}(f)$ & $\diff{N^{+}}(f)$ & Notation & \textbf{Umgesetzt.} Subdifferential uneinheitlich: $\ensuremath{\partial}N\ensuremath{^{+}}(f)$ in Z. 61, 76, 95, 215, dagegen $dN\ensuremath{^{+}}(f)$ mit aufrechtem d \code{(\allowbreak{}\textbackslash{}diff)} in Z. 91, 106, 181, 227, 250, 253, 487 und \code{\textbackslash{}ddp(f)} in Z. 1094, 1106, 1596. Teil I schreibt \code{\textbackslash{}ddp}, daher Vereinheitlichung auf d vorgeschlagen; ebenso Z. 76, 95, 215. \\[2pt]
280 & 64 & \sichB{} & {see Eq.~(\lbl{eq:a2-2.5}).} & {see Eq.~(\lbl{eq:a2-2.5})).} & Klammern & \textbf{Umgesetzt.} Schon in der Arbeitsliste. Die mit „(this follows \dots{}“ geöffnete Klammer wird nicht geschlossen. \\[2pt]
281 & 190 & \sichB{} & $C^{2}[0,1]$ & $C^{2}([0,1])$ & Notation & \textbf{Umgesetzt.} Sonst im Kapitel $C^{2}([0,1])$, $L^{p}([0,1])$, AC([0,1]); ebenso Z. 1287 ($L^{q}[0,1]$) und Z. 2074, 2123, 2136 ($C[-1,0]$, $C^{1}[-1,0]$ gegenüber $C^{1}([-1,0],\C)$ in Z. 2046, 2055). \\[2pt]
281 & 216 & \sichA{} & {\lbl{ex:c2-1.1},~\lbl{ex:c2-1.1-2} )} & {\lbl{ex:c2-1.1}\,\lbl{ex:c2-1.1-2})} & Abstand & \textbf{Umgesetzt.} Leerzeichen vor der schließenden Klammer; außerdem sonst „Example~\lbl{…}\textbackslash{},\lbl{…}“ ohne Komma (Z. 1786, 1792, 2243). \\[2pt]
281 & 252 & \sichB{} & $f(x) = \|f\|_{\infty}$ & $f(x) = \|f^{+}\|_{\infty}$ & Tippfehler & \textbf{Umgesetzt.} bitte prüfen: Das Maximum von f ist $\ensuremath{\|}f\ensuremath{^{+}}\ensuremath{\|}\ensuremath{\infty}$ (vgl. Z. 75, 179), nicht $\ensuremath{\|}f\ensuremath{\|}\ensuremath{\infty}$. \\[2pt]
283 & 264 & \sichA{} & $(1 + \sum_{i=1}^{n} y_{i}^{2})g(y), g \in \mathcal{S}(\R^{n})) .$ & $(1 + \sum_{i=1}^{n} y_{i}^{2})g(y) \quad (g \in \mathcal{S}(\R^{n})) .$ & Klammern & \textbf{Umgesetzt.} Überzählige schließende Klammer; die Nebenbedingung gehört offenbar in Klammern. \\[2pt]
283 & 308 & \sichB{} & {\,$C_{c}^{\infty}(\R^{n})$\, is dense in \,$\mathcal{S}$\,} & {$C_{c}^{\infty}(\R^{n})$ is dense in $\mathcal{S}$} & Abstand & \textbf{Umgesetzt.} Die manuellen \code{\textbackslash{}}, um die Inline-Formeln vergrößern die Wortabstände sichtbar. \\[2pt]
283 & 317 & \sichA{} & $\phi \in S$ & $\phi \in \mathcal{S}$ & Notation & \textbf{Umgesetzt.} Der Schwartz-Raum heißt sonst $\mathcal{S}$ (Z. 286, 314). \\[2pt]
284 & 340 & \sichC{} & $u(x) - \Delta u(x)  = 0 \quad  (x \in \Omega)$ & \code{\textbackslash{}begin\{array\}\{r@\{\allowbreak{}\textbackslash{};\}l@\{\allowbreak{}\textbackslash{}quad\}l\}} & Abstand & \textbf{Umgesetzt.} Mit {aligned} gelöst (ebenso Z. 1286). Im array steht vor dem „$=$“ der Spaltenabstand \code{2\allowbreak{}\textbackslash{}arraycolsep} (10 pt), dahinter nur der normale Relationsabstand. Alternativ aligned mit \&=. Ebenso Z. 1286 (\{rl\}, $D(A')  = \dots$). \\[2pt]
284 & 381 & \sichB{} & $f \in C^{2}(\Omega) \cap C(\Omega)$ & $f \in C^{2}(\Omega) \cap C(\overline{\Omega})$ & Tippfehler & \textbf{Umgesetzt.} bitte prüfen: $C(\ensuremath{\Omega})$ ist in $C\ensuremath{^{2}}(\ensuremath{\Omega})$ bereits enthalten; aus \ensuremath{\bar{f}} $\ensuremath{\in} \ensuremath{\mathcal{S}}$ und $u \ensuremath{\in} C(\ensuremath{\bar{\ensuremath{\Omega}}})$ folgt $f \ensuremath{\in} C(\ensuremath{\bar{\ensuremath{\Omega}}})$ (vgl. Z. 376). \\[2pt]
284 & 399 & \sichB{} & {We know from~\lbl{ex:c2-1.5-4}} & {We know from~\lbl{ex:c2-1.5-5}} & Verweis & \textbf{Umgesetzt.} Schon in der Arbeitsliste. Die Aussage wird im Beweis von (v) gezeigt. \\[2pt]
285 & 426 & \sichA{} & {$E = C(K)$ (K compact)} & {$E = C(K)$ ($K$ compact)} & Mathemodus & \textbf{Umgesetzt.} K steht nicht im Mathemodus; ebenso Z. 2024 ({(X locally compact)}), Z. 2399 ({Assume that E is}), Z. 2885 ({(K compact)}). \\[2pt]
286 & 518 & \sichB{} & $\lim_{n \to \infty} \|u - (n u \wedge g)\| = 0$ & $\lim_{n \to \infty} \|g - (n u \wedge g)\| = 0$ & Tippfehler & \textbf{Umgesetzt.} bitte prüfen: g liegt genau dann im Abschluss von $E_u$, wenn $nu \ensuremath{\wedge} g$ gegen g konvergiert. \\[2pt]
286 & 525 & \sichB{} & $(Au)^{-} \in E_{u}$ & $(Au)^{-} \in \overline{E_{u}}$ & Tippfehler & \textbf{Umgesetzt.} bitte prüfen: Der Beweis (Z. 532, 536) und Theorem 2.1.11 (c) (Z. 571) verwenden den Abschluss. \\[2pt]
286 & 532 & \sichC{} & $\overline{E}_{u}$ & $\overline{E_{u}}$ & Notation & \textbf{Umgesetzt.} Z. 532, 533, 536 (3 Stellen). Querstrich hier nur über E, sonst über $E_u$ bzw. $E_f$ (Z. 518, 571, 1378); ebenso Z. 533 $(2\ensuremath{\times})$ und 536. \\[2pt]
286 & 533 & \sichC{} & $0 \leq \phi \in E_{+}'$ & $0 \leq \phi \in E'_{+}$ & Notation & \textbf{Umgesetzt.} Sonst durchgehend $E'_{+}$; im Druck steht der Strich hier hinter dem Index. \\[2pt]
287 & 571 & \sichC{} & \code{\textbackslash{}overline\{E\_\{f\}\}\$ \allowbreak{}\allowbreak{}\textbackslash{} \allowbreak{}for \allowbreak{}all} & \code{\textbackslash{}overline\{E\_\{f\}\}\$ \allowbreak{}for \allowbreak{}all} & Abstand & \textbf{Umgesetzt.} Leerzeichen plus „\code{\textbackslash{}} “ ergibt doppelten Wortabstand; ebenso Z. 1848 ($\text{for }\ 0$). \\[2pt]
287 & 610 & \sichB{} & $Pg = g^{-} < 0$ & $Pg = -g^{-} < 0$ & Tippfehler & \textbf{Umgesetzt.} bitte prüfen: $g\ensuremath{^{-}} \ensuremath{\geq} 0$; P projiziert auf das von $g\ensuremath{^{-}}$ erzeugte Band, also $Pg = \ensuremath{-}g\ensuremath{^{-}}$. \\[2pt]
288 & 653 & \sichA{} & {the $n^{th}$ unit vector} & {the $n$th unit vector} & Mathemodus & \textbf{Umgesetzt.} „th“ wird als kursives Variablenprodukt im Exponenten gesetzt. \\[2pt]
288 & 671 & \sichB{} & {linear operator (sign $f$) on $E$} & {linear operator $(\sign f)$ on $E$} & Notation & \textbf{Umgesetzt.} Im kursiven Satztext erscheint „sign“ kursiv, sonst überall aufrecht als \code{\textbackslash{}sign}. \\[2pt]
289 & 720 & \sichB{} & $(\sign f)g = m \cdot g \quad (g \in E),$ & $(\sign f)g = m \cdot g \quad (g \in E).$ & Satzzeichen & \textbf{Umgesetzt.} Nach der Formel beginnt ein neuer Satz („Moreover,“). \\[2pt]
289 & 750 & \sichC{} & $\Re\left(\widehat{\sign} \bar{f}\right)g$ & $\Re\left((\widehat{\sign} \bar{f})g\right)$ & Klammern & \textbf{Umgesetzt.} So steht der Realteil des Operators statt des Elements; in Z. 2450 dieselbe Formel als \code{\textbackslash{}Re((ŝign} \ensuremath{\bar{f}})g). Ebenso Z. 2491, 2526, 2564, 2566. \\[2pt]
290 & 799 & \sichC{} & \code{\textbackslash{}begin\{equation*\}\allowbreak{}\textbackslash{}label\{eq:c2-K-2\}} & \code{\textbackslash{}begin\{equation*\}} & LaTeX & \textbf{Umgesetzt.} \code{\textbackslash{}label} in unnummerierter Formel ohne \code{\textbackslash{}tag}, nirgends referenziert; ebenso Z. 1355 (eq:c2-3-KGK). Im Druck unsichtbar, nur Bereinigung. \\[2pt]
291 & 848 & \sichA{} & $\langle \phi,\Delta\vert f\vert \rangle (= \langle \Delta \phi,\vert f\vert \rangle$ & $\langle \phi,\Delta\vert f\vert \rangle \;(= \langle \Delta \phi,\vert f\vert \rangle)$ & Klammern & \textbf{Umgesetzt.} Die mit „($=$“ geöffnete Klammer wird nicht geschlossen. \\[2pt]
291 & 851 & \sichA{} & $(1~\leq~p <~\infty)$ & $(1 \leq p < \infty)$ & Abstand & \textbf{Umgesetzt.} Tilden im Mathemodus erzeugen zusätzliche Wortabstände; ebenso Z. 1412 ($-\phi^{(3)}(x)~\leq~\phi(x)$) und Z. 1422 ($f~\in~D(B)$). \\[2pt]
291 & 884 & \sichC{} & $C_{c}^{\infty} \subseteq D_{0}$ & $C_{c}^{\infty}(\R^n) \subseteq D_{0}$ & Notation & \textbf{Umgesetzt.} Im übrigen Absatz $C_c^\ensuremath{\infty}(R^n)$; die Kurzform $C_c^\ensuremath{\infty}$ wird erst in Z. 1792 eingeführt. \\[2pt]
291 & 893 & \sichA{} & $\Re\langle(\sign \bar{f})Af),\phi \rangle$ & $\Re\langle(\sign \bar{f})Af,\phi \rangle$ & Klammern & \textbf{Umgesetzt.} Überzählige „)“. \\[2pt]
291 & 894 & \sichB{} & $\leq \langle \vert f\vert,A'\phi \rangle   = \langle$ & $\leq \langle \vert f\vert,A'\phi \rangle = \langle$ & LaTeX & \textbf{Umgesetzt.} Das \& der zweiten Zeile steht an der falschen Stelle; das „$\ensuremath{\leq}$“ wird nicht unter dem „$=$“ der ersten Zeile ausgerichtet. \\[2pt]
292 & 901 & \sichB{} & $\sum_{i,j=1}^{2} b_{ij}$ & $\sum_{i,j=1}^{n} b_{ij}$ & Tippfehler & \textbf{Umgesetzt.} bitte prüfen: Der Operator wirkt auf dem $R^n$; die obere Grenze 2 ist vermutlich mit der Ordnung $m \ensuremath{\leq} 2$ verwechselt. \\[2pt]
294 & 1032 & \sichB{} & {Definition~\lbl{def:b3-3.1}} & {Definition~\lbl{def:c3-3.1}} & Verweis & \textbf{Umgesetzt.} bitte prüfen: def:b3-3.1 ist die Irreduzibilität auf $C_0(X)$ in Teil II (hier ohne „Part-II,“); für Banachverbände gibt es Definition 3.3.1 in Teil III (def:c3-3.1). Gedruckte Nummer gleich, Link und Seitenbezug führen aber nach Teil II. \\[2pt]
294 & 1053 & \sichB{} & $f \in D_{0} \, ,  \phi \in M'$ & $f \in D_{0}, \phi \in M'$ & Abstand & \textbf{Umgesetzt.} Sichtbarer Abstand vor dem Komma (vgl. die gleiche Formel in Z. 1205). \\[2pt]
295 & 1087 & \sichC{} & $P_{+} \coloneq P_{f^{+}}$ & $P_{+} \coloneq P_{(f^{+})}$ & Notation & \textbf{Umgesetzt.} Bandprojektion teils $P_{(f^{+})}$ (Z. 700, 1182, 1187, 1740, 1823), teils $P_{f^{+}}$ (Z. 1087, 1932, 2658); einheitlich setzen. \\[2pt]
295 & 1089 & \sichB{} & $Q \coloneq P_{+} + 1/2$ & $Q \coloneq P_{+} + 1/2\,P$ & Tippfehler & \textbf{Umgesetzt.} bitte prüfen: Nach Z. 1104 ist $Qf = P\ensuremath{_{+}}f + \ensuremath{\tfrac{1}{2}}Pf$; „$+ 1/2$“ allein ergäbe auch nicht $Q \ensuremath{\leq} \Id $. \\[2pt]
295 & 1139 & \sichA{} & $\text{ and } t\geq 0) .$ & $\text{ and } t\geq 0 .$ & Klammern & \textbf{Umgesetzt.} Überzählige „)“. \\[2pt]
297 & 1260 & \sichA{} & $f(1) = \diff{}f(0)$ & $f(1) = d f(0)$ & Tippfehler & \textbf{Umgesetzt.} \code{\textbackslash{}diff\{\}} setzt das aufrechte Differential-d; gemeint ist die Konstante d $\ensuremath{\in} R$ aus Z. 1256 (vgl. $\ensuremath{\varphi}(0) = d\ensuremath{\varphi}(1)$ in Z. 1287). \\[2pt]
297 & 1266 & \sichC{} & $(n \in \N )$ & $(n \in \N_{0})$ & Tippfehler & \textbf{Umgesetzt.} bitte prüfen: Für $0 \ensuremath{\leq} x+t < 1$ wird $n = 0$ gebraucht; $\N_0$ ist in Z. 872 eingeführt. \\[2pt]
298 & 1288 & \sichB{} & $\text{for all } \phi \in D(A').$ & $\text{for all } \phi \in D(A'),$ & Satzzeichen & \textbf{Umgesetzt.} Nach dem Punkt geht der Satz mit „where $1/p + 1/q = 1$“ weiter. \\[2pt]
298 & 1324 & \sichC{} & $\text{has codimension one in } D(A_{\max})$ & $\text{has codimension one in } D(A_{\max}).$ & Satzzeichen & \textbf{Umgesetzt.} Die übrigen Zeilen dieser Aufzählung enden mit Punkt. \\[2pt]
299 & 1366 & \sichA{} & $\phi \in D(B')_{+}) .$ & $\phi \in D(B')_{+} .$ & Klammern & \textbf{Umgesetzt.} Überzählige „)“. \\[2pt]
300 & 1439 & \sichA{} & $\left(-\phi^{(3)}(x)\right)) \dx$ & $\left(-\phi^{(3)}(x)\right) \dx$ & Klammern & \textbf{Umgesetzt.} Überzählige „)“. \\[2pt]
301 & 1511 & \sichA{} & {Section~\lbl{sec:c3-1}} & {Section~\lbl{sec:c2-3}} & Verweis & \textbf{Umgesetzt.} Schon in der Arbeitsliste. Die Charakterisierung positiver Halbgruppen steht in Abschnitt 2.3, nicht in 3.1. \\[2pt]
301 & 1531 & \sichB{} & $\Re\mleft\langle(\sign  \bar{f}) S(t)f\mright\rangle$ & $\Re\mleft[(\sign \bar{f}) S(t)f\mright]$ & Klammern & \textbf{Umgesetzt.} bitte prüfen: Spitze Klammern mit nur einem Argument; gemeint ist der Realteil des Elements (vgl. Z. 1719 \code{\textbackslash{}Re[\dots{}]}). \\[2pt]
301 & 1537 & \sichC{} & $\max\{\omega (A),\omega (B)\}$ & $\max\{\omega_{0}(A),\omega_{0}(B)\}$ & Notation & \textbf{Umgesetzt.} Wachstumsschranke sonst $\ensuremath{\omega}_0$ (Z. 407, 1677, 2218). \\[2pt]
302 & 1552 & \sichC{} & $= \langle R(\lambda,A)\vert g\vert,\psi\rangle$ & $= \langle R(\lambda,A)\vert g\vert,\psi\rangle .$ & Satzzeichen & \textbf{Umgesetzt.} Abgesetzte Formel ohne Satzende, danach beginnt ein neuer Satz; ebenso Z. 1600 (vor „Thus“) und Z. 1642 (vor „Assume“). \\[2pt]
303 & 1651 & \sichB{} & $\vert f_{n} - f_{m}\vert \leq R(\lambda,A)\vert g_{n} - g_{m}\vert$ & $\vert f_{n} - f_{m}\vert \leq R(\mu,A)\vert g_{n} - g_{m}\vert$ & Tippfehler & \textbf{Umgesetzt.} bitte prüfen: $g_n =$ ($\ensuremath{\mu} \ensuremath{-} B)f_n$, also ist $R(\ensuremath{\mu},A)$ gemeint. \\[2pt]
304 & 1694 & \sichB{} & {(the center of\/ $E$ (see Section~\lbl{sec:c1-9}).} & {(the center of\/ $E$, see Section~\lbl{sec:c1-9}).} & Klammern & \textbf{Umgesetzt.} Schon in der Arbeitsliste. Klammern nicht ausgeglichen. \\[2pt]
304 & 1695 & \sichC{} & {$E =$ $L^{p}(X,\mu)$} & $E = L^{p}(X,\mu)$ & LaTeX & \textbf{Umgesetzt.} Formel unnötig geteilt (Umbruch nach „$=$“ möglich, Abstand nach „$=$“ ist Wortabstand); ebenso Z. 2218 ({$\lambda \geq$ $\omega_{0}(A)+1$}). \\[2pt]
304 & 1711 & \sichB{} & $B_{0} \in \mathcal{L}(E^{r})$ & $B_{0} \in \mathcal{L}^{r}(E)$ & Notation & \textbf{Umgesetzt.} bitte prüfen: Reguläre Operatoren sonst $\mathcal{L}^{r}(E)$ (Z. 2207, 2280, 2294); in Z. 2207 genauer $B_0 \ensuremath{\in} \mathcal{Z}(E)^{d}$. \\[2pt]
307 & 1924 & \sichB{} & $\lambda R(\lambda,A)u = \lambda R(\lambda,A)u+u$ & $\lambda R(\lambda,A)u = AR(\lambda,A)u+u$ & Tippfehler & \textbf{Umgesetzt.} bitte prüfen: So unterscheiden sich beide Seiten um u; gemeint ist $\ensuremath{\lambda}R(\ensuremath{\lambda},A)u = AR(\ensuremath{\lambda},A)u + u$. \\[2pt]
309 & 2071 & \sichC{} & $\alpha \coloneq \Re(\beta)$ & $\alpha = \Re(\beta)$ & Notation & \textbf{Umgesetzt.} $\ensuremath{\alpha}$ ist vorgegeben; formuliert wird eine Annahme, keine Definition (parallel dazu „$\mu = \vert\nu\vert$“). \\[2pt]
309 & 2086 & \sichA{} & $\leq \langle\vert f\vert,A'\phi\rangle .$ & $\leq \langle\vert f\vert,A'\phi\rangle$ & Satzzeichen & \textbf{Umgesetzt.} Punkt am Formelende, obwohl der Satz mit „holds for all \dots{}“ weitergeht. \\[2pt]
310 & 2129 & \sichA{} & $1/t \, ((T(t)\vert f\vert)(0)$ & $1/t \, (T(t)\vert f\vert)(0)$ & Klammern & \textbf{Umgesetzt.} Eine öffnende Klammer zu viel. \\[2pt]
310 & 2134 & \sichC{} & $\mu\{0\} = \nu\{0\} = 0$ & $\mu(\{0\}) = \nu(\{0\}) = 0$ & Notation & \textbf{Umgesetzt.} In Z. 2041 $\ensuremath{\mu}({0})$; ebenso Z. 2144. \\[2pt]
311 & 2181 & \sichB{} & $- \langle\Re(V)\vert f\vert,\phi\rangle  = \Re\langle$ & $- \langle\Re(V)\vert f\vert,\phi\rangle = \Re\langle$ & LaTeX & \textbf{Umgesetzt.} Das zweite \& öffnet eine neue align-Spaltengruppe; „$= \Re \ensuremath{\langle}\dots Af\dots \ensuremath{\rangle}$“ rückt mit großem Abstand nach rechts. \\[2pt]
312 & 2234 & \sichC{} & $Mc/\lambda(\|g^{+}\| + \|g^{-}\|)$ & $(Mc/\lambda)(\|g^{+}\| + \|g^{-}\|)$ & Klammern & \textbf{Umgesetzt.} Ohne Klammern liest sich $\ensuremath{\lambda}(\dots )$ als Nenner; vgl. Z. 2229 und rechts $(2Mc/\ensuremath{\lambda})$. \\[2pt]
314 & 2354 & \sichB{} & {Then by Eq.~(\lbl{eq:c2-5.1}),} & {Then by Eq.~(\lbl{eq:c2-5.3}),} & Verweis & \textbf{Umgesetzt.} bitte prüfen: Benutzt wird die Definition von T(s) als Modulus von S(s), also (2.5.3), nicht die Verbandshomomorphismus-Gleichung (2.5.1). \\[2pt]
315 & 2489 & \sichC{} & $\displaystyle D_{g}\Theta(f) = \lim_{t \downarrow 0} k(t)/t$ & $D_{g}\Theta(f) = \lim_{t \downarrow 0} k(t)/t$ & LaTeX & \textbf{Umgesetzt.} \code{\textbackslash{}displaystyle} im Fließtext vergrößert den Zeilenabstand; ggf. als abgesetzte Formel. \\[2pt]
316 & 2505 & \sichB{} & {Part-II, Chapter~\lbl{chap:b2}, \ensuremath{\hookleftarrow} Lemma 2.4} & {Part-II, Lemma~\lbl{lem:b2-2.4}} & Verweis & \textbf{Umgesetzt.} Schon in der Arbeitsliste. Hartkodierte Nummer und abweichende Zitierweise (Z. 737: Part-II, Lemma~\lbl{lem:b2-2.4}). \\[2pt]
316 & 2512 & \sichC{} & {regular Borel measure $\mu$ on $C(K)$} & {regular Borel measure $\mu$ on $K$} & Tippfehler & \textbf{Umgesetzt.} bitte prüfen: Das Maß lebt auf K (integriert wird in Z. 2517 über K). \\[2pt]
316 & 2517 & \sichA{} & $\Re\mleft(\sign (\bar{f}(x))\mright)h(x))$ & $\Re\mleft(\sign(\overline{f(x)})h(x)\mright)$ & Klammern & \textbf{Umgesetzt.} Klammern unausgeglichen: \code{\textbackslash{}mright}) steht zu früh, am Ende eine „)“ zu viel; vgl. Z. 2509. \\[2pt]
316 & 2524 & \sichC{} & {(see Section~\lbl{sec:c1-8})\,).} & {(see Section~\lbl{sec:c1-8})).} & Abstand & \textbf{Umgesetzt.} Schon in der Arbeitsliste. Abstand vor der schließenden Klammer. \\[2pt]
317 & 2539 & \sichC{} & $\Re((\sign  \overline{f})g)$ & $\Re((\sign  \bar{f})g)$ & Notation & \textbf{Umgesetzt.} Sonst durchgehend \code{\textbackslash{}bar\{f\}}; \code{\textbackslash{}overline\{f\}} ist im Druck breiter. \\[2pt]
317 & 2542 & \sichB{} & $D^{-}_{g}\Theta(f) = \lim_{t \to 0}$ & $D^{-}_{g}\Theta(f) = \lim_{t \uparrow 0}$ & Tippfehler & \textbf{Umgesetzt.} bitte prüfen: Linksseitige Ableitung; die rechtsseitige ist mit $t \ensuremath{\downarrow} 0$ notiert (Z. 2489). \\[2pt]
317 & 2578 & \sichA{} & $\text{ for } t \in [0,s]) .$ & $\text{ for } t \in [0,s] .$ & Klammern & \textbf{Umgesetzt.} War im Quelltext schon korrekt. Überzählige „)“. \\[2pt]
318 & 2664 & \sichB{} & $(P_{+} - P_{-})(Af^{+} + Af^{-})$ & $(P_{+} - P_{-})(Af^{+} - Af^{-})$ & Tippfehler & \textbf{Umgesetzt.} bitte prüfen: $Af = Af\ensuremath{^{+}} \ensuremath{-} Af\ensuremath{^{-}}$; die folgende Zeile setzt das Minuszeichen voraus. \\[2pt]
320 & 2747 & \sichC{} & {is a multiplication operator for all $\lambda \in \rho(A)$} & {is a multiplication operator for all $\lambda \in \rho(A)$.} & Satzzeichen & \textbf{Umgesetzt.} In der Aufzählung enden (a), (b), (e) mit Punkt, (c) und (d) (Z. 2750) nicht; ebenso Prop. 2.5.15 (a), (b) (Z. 2833, 2836). \\[2pt]
321 & 2846 & \sichA{} & {\lbl{sec:c1-9})\ensuremath{\hookleftarrow}, it is clear} & {\lbl{sec:c1-9}), it is clear} & Abstand & \textbf{Umgesetzt.} Schon in der Arbeitsliste. Der Zeilenumbruch vor dem Komma ergibt im Druck ein Leerzeichen. \\[2pt]
321 & 2849 & \sichB{} & {Hence $A \in \mathcal{Z}(E)$} & {Hence $A \in \mathcal{Z}(E)$.} & Satzzeichen & \textbf{Umgesetzt.} Schon in der Arbeitsliste. Am Satzende fehlt der Punkt. \\[2pt]
322 & 2898 & \sichA{} & $\phi \in E^*$ & $\phi \in E'$ & Notation & \textbf{Umgesetzt.} Dualraum sonst durchgehend E'; ebenso Z. 2922 ($E_{+}^*$ $\ensuremath{\to}$ $E'_{+}$). \\[2pt]
322 & 2899 & \sichB{} & $Af \in (f)^{dd}$ & $Af \in \{f\}^{dd}$ & Notation & \textbf{Umgesetzt.} Sonst $\{f\}^{dd}$ (Z. 1378, 2492, 2633). \\[2pt]
322 & 2906 & \sichB{} & {(this follows from Eq.~(\lbl{eq:a2-1.8}))} & {(this follows from Part-I, Eq.~(\lbl{eq:a2-1.8}))} & Verweis & \textbf{Umgesetzt.} Schon in der Arbeitsliste. Formel aus Teil I ohne Teilangabe; die gedruckte Nummer kommt auch in diesem Kapitel vor. \\[2pt]
322 & 2944 & \sichC{} & {C*-algebras} & {\CA-algebras} & Notation & \textbf{Umgesetzt.} In Z. 2975 mit {\CA{}} (C mit hochgestelltem Stern), hier mit Text-Sternchen. \\[2pt]
322 & 2951 & \sichC{} & $C^{\infty}_0(\R^n)$ & $C^{\infty}_{c}(\R^n)$ & Notation & \textbf{Umgesetzt.} Testfunktionen sonst $C_c^\ensuremath{\infty}$ (z. B. Z. 1785, 2166). \\
\end{longtable}
\subsection*{Teil III, Kap. 3 – Spectral Theory of Positive Semigroups on Banach Lattices (\code{part-c/chap-c3.tex})}
\begin{longtable}{@{}rrc>{\RaggedRight\arraybackslash}p{5.8cm}>{\RaggedRight\arraybackslash}p{5.8cm}>{\RaggedRight\arraybackslash}p{1.9cm}>{\RaggedRight\arraybackslash}p{7.3cm}@{}}
\toprule
\textbf{S.} & \textbf{Z.} & \textbf{} & \textbf{Ist} & \textbf{Soll} & \textbf{Art} & \textbf{Kommentar} \\
\midrule
\endfirsthead
\toprule
\textbf{S.} & \textbf{Z.} & \textbf{} & \textbf{Ist} & \textbf{Soll} & \textbf{Art} & \textbf{Kommentar} \\
\midrule
\endhead
\bottomrule
\endlastfoot
327 & 52 & \sichC{} & {If \ $T(1)'\phi_{0} = \phi_{0}$ \ for some \ $\phi_{0} \in E'_{+}$} & {If $T(1)'\phi_{0} = \phi_{0}$ for some $\phi_{0} \in E'_{+}$} & Abstand & Erzwungene Leerzeichen „\code{\textbackslash{}} “ zusätzlich zum normalen Wortabstand ergeben sichtbar doppelte Lücken. \\[2pt]
328 & 115 & \sichC{} & \code{T(s)\allowbreak{}\textbackslash{}vert \allowbreak{}f\allowbreak{}\textbackslash{}vert \allowbreak{}\allowbreak{}\textbackslash{}ds \allowbreak{}\allowbreak{}\textbackslash{}right\allowbreak{}\textbackslash{}| \allowbreak{}\allowbreak{}\textbackslash{}} & \code{T(s)\allowbreak{}\textbackslash{}vert \allowbreak{}f\allowbreak{}\textbackslash{}vert \allowbreak{}\allowbreak{}\textbackslash{}ds \allowbreak{}\allowbreak{}\textbackslash{}right\allowbreak{}\textbackslash{}| \allowbreak{}.} & LaTeX & Überzähliger Backslash am Zeilenende (wirkt als Leerzeichen), Satzzeichen fehlt; ebenso Z. 684 ($= 0\ $). \\[2pt]
328 & 121 & \sichB{} & $\left(\int_{0}^t \exp(-\mu s)T(s) \ds \right)_{t\geq 0}$ & $\left(\int_{0}^t \eu^{-\lambda s}T(s) \ds \right)_{t\geq 0}$ & Tippfehler & Bitte prüfen: Zu zeigen ist die Normkonvergenz für $\ensuremath{\lambda}$ (vgl. Aussage Z. 76 und Schluss Z. 138); $\ensuremath{\mu}$ wird erst in Z. 122 fixiert. \\[2pt]
329 & 134 & \sichB{} & $\exp(-(\Re(\lambda)-\mu)s)\eu^{-\lambda s}\vert\langle T(s)f,\phi\rangle\vert$ & $\exp(-(\Re(\lambda)-\mu)s)\eu^{-\mu s}\vert\langle T(s)f,\phi\rangle\vert$ & Tippfehler & Bitte prüfen: Für die Abschätzung mit Eq. (3.1.4) muss der Faktor $e^{-\ensuremath{\mu}s}$ stehen; ebenso Z. 135 ($\int_{t}^r \eu^{-\lambda s}\vert\langle$ $\ensuremath{\to}$ $\eu^{-\mu s}$). \\[2pt]
331 & 223 & \sichB{} & \code{\textbackslash{}index\{spectrum!boundary\allowbreak{}\textbackslash{}vert)\} \allowbreak{}\dots{} \allowbreak{}\allowbreak{}\textbackslash{}section\{The \allowbreak{}Boundary \allowbreak{}Spectrum\}} & \code{\textbackslash{}index\{spectrum!boundary\allowbreak{}\textbackslash{}vert)\} \allowbreak{}\allowbreak{}\textbackslash{}section\{Irreducible \allowbreak{}Semigroups\}} & LaTeX & Kein Formelbefund: Indexbereich wird in Z. 223 geschlossen, aber erst in Z. 227 geöffnet und danach nie geschlossen (makeindex-Warnung, falscher Seitenbereich). \\[2pt]
331 & 262 & \sichC{} & $\vert f^{[n]}\vert = \vert f\vert \quad (n \in \Z)$ & $\vert f^{[n]}\vert = \vert f\vert \quad (n \in \Z);$ & Satzzeichen & In der Aufzählung Eq. (3.2.3)–(3.2.5) endet (3.2.4) mit „;“, (3.2.5) mit „.“; bei (3.2.3) fehlt das Semikolon. \\[2pt]
332 & 305 & \sichB{} & $\overline{T}(t)\overline{h}^{[n]} = \eu^{\im\beta  t}\overline{h}^{[n]}$ & $\overline{T}(t)\overline{h}^{[n]} = \eu^{\im n\beta t}\overline{h}^{[n]}$ & Tippfehler & Bitte prüfen: Es fehlt der Faktor n im Exponenten (vgl. Z. 306 und Eq. (3.2.7)). \\[2pt]
332 & 306 & \sichB{} & $T(t)h^{[n]} = \eu^{\im n\beta}h^{[n]}$ & $T(t)h^{[n]} = \eu^{\im n\beta t}h^{[n]}$ & Tippfehler & Bitte prüfen: Im Exponenten fehlt t. \\[2pt]
333 & 373 & \sichC{} & $0 \in P\sigma(R_{\F}(\lambda,A))$ & $0 \in P\sigma(R_{\F}(\lambda))$ & Notation & Die Pseudoresolvente $R_F$ hat ein Argument (Z. 372: „$R(\ensuremath{\lambda},A)_F$“ bzw. „$R_F$“); bitte prüfen. \\[2pt]
333 & 397 & \sichC{} & \code{R(\allowbreak{}\textbackslash{}lambda\_\{0\})\textasciicircum{}\{n+1\}\$ \allowbreak{}\} \allowbreak{}.} & \code{R(\allowbreak{}\textbackslash{}lambda\_\{0\})\textasciicircum{}\{n+1\}\$\} \allowbreak{}.} & Abstand & Das Leerzeichen vor „\}“ steht innerhalb von \code{\textbackslash{}text} und ist im Druck als Lücke vor dem Punkt sichtbar. \\[2pt]
334 & 420 & \sichC{} & $S = (\mu-\lambda_{0})^{-1}(1 - (1-(\lambda_{0}-\mu)R(\lambda_{0}))^{-1})$ & $S = (\mu-\lambda_{0})^{-1}(\Id - (\Id-(\lambda_{0}-\mu)R(\lambda_{0}))^{-1})$ & Notation & Unmittelbar davor (Z. 413 f.) steht \code{\textbackslash{}Id} für dieselbe Identität; ebenso Z. 1635 (vgl. Z. 1627/1630 mit \code{\textbackslash{}Id}) und Z. 1680; außerdem 1 statt \code{\textbackslash{}Id} in Z. 357, 658, 661. \\[2pt]
334 & 435 & \sichB{} & $\overline{D}\backslash D$ & $\overline{D}\setminus D$ & Notation & \code{\textbackslash{}backslash} ist sichtbar größer als \code{\textbackslash{}setminus \allowbreak{}(\allowbreak{}\textbackslash{}smallsetminus)}, das sonst im Kapitel verwendet wird (Z. 1085, 1166, 1336); ebenso Z. 1775, 1902. \\[2pt]
335 & 508 & \sichB{} & $S_{h}^{n}\vert h\vert = S_{h}^{-n}z\tilde{R}(z) \vert h\vert$ & $S_{h}^{n}\vert h\vert = S_{h}^{n}z\tilde{R}(z) \vert h\vert$ & Tippfehler & Bitte prüfen: Nach Eq. (3.2.14) muss links $S_h^{n}$ stehen, sonst passt der Folgeschritt nicht. \\[2pt]
336 & 547 & \sichC{} & {Eq.~(\lbl{eq:c3-2.15})~(a)} & {condition (a) of Eq.~(\lbl{eq:c3-2.15})} & Verweis & \code{\textbackslash{}myeqpageref} hängt ggf. „on the \dots{} page“ an, dann steht „(a)“ hinter der Seitenangabe; ebenso Z. 554 und 564. \\[2pt]
336 & 602 & \sichB{} & $r\lambda \eu^{-r}$ & $r\lambda \eu^{-\lambda r}$ & Tippfehler & Bitte prüfen: Aus Z. 598 folgt der Faktor $e^{-\ensuremath{\lambda}r}$, nur dann ergibt sich ($1 - e^{-\ensuremath{\lambda}r}$) in Z. 603. \\[2pt]
336 & 606 & \sichB{} & {hence $\|\lambda R(\lambda,A)\| \leq M$} & {hence $\|\lambda R(\lambda,A)\| \leq M$.} & Satzzeichen & Der Satz am Beweisende hat keinen Punkt. \\[2pt]
337 & 625 & \sichB{} & {in Corollary~\lbl{cor:c3-2.12}} & {in Corollary~\lbl{cor:c3-2.12}.} & Satzzeichen & Nach dem Verweis fehlt der Schlusspunkt. \\[2pt]
338 & 655 & \sichB{} & {for $r > 0$, $\epsilon > 0$} & {for $\tau > 0$, $\epsilon > 0$} & Tippfehler & Die folgende Formel und Z. 661 verwenden $\ensuremath{\tau}$, nicht r. \\[2pt]
338 & 688 & \sichA{} & $r > 0) .$ & $r > 0 .$ & Klammern & Schließende Klammer ohne öffnende Klammer. \\[2pt]
338 & 690 & \sichA{} & $\{(R_{1}(\lambda)\}_{\Re(\lambda)>0}$ & $\{R_{1}(\lambda)\}_{\Re(\lambda)>0}$ & Klammern & Öffnende Klammer ohne schließende Klammer. \\[2pt]
338 & 691 & \sichA{} & {(by Eq.~(\lbl{eq:c3-2.19}).} & {(by Eq.~(\lbl{eq:c3-2.19})).} & Klammern & Schon in der Arbeitsliste. Ebenso Z. 692: {(by Eq.~(\lbl{eq:c3-2.16}) and} $\ensuremath{\to}$ {(by Eq.~(\lbl{eq:c3-2.16})) and}. \\[2pt]
338 & 694 & \sichB{} & {Proposition~\lbl{prop:c3-2.7}(b)} & {Proposition~\lbl{prop:c3-2.7}} & Verweis & Schon in der Arbeitsliste. Proposition 3.2.7 hat keinen Teil (b). \\[2pt]
339 & 716 & \sichC{} & $\subsetneq ... \subsetneq$ & $\subsetneq \dots \subsetneq$ & LaTeX & „...“ im Mathemodus ergibt schlecht gesetzte Punkte; sonst \code{\textbackslash{}ldots}. \\[2pt]
339 & 722 & \sichB{} & {$s(A_{0}) < s(A)$ ,} & {$s(A_{0}) < s(A)$,} & Abstand & Sichtbares Leerzeichen vor dem Komma; ebenso Z. 2056 (\code{(R(\allowbreak{}\textbackslash{}mu,A)g)\textasciicircum{}+\$ \allowbreak{},}) und Z. 2071 (\code{R(\allowbreak{}\textbackslash{}mu,A)h \allowbreak{}\allowbreak{}\textbackslash{}geq \allowbreak{}0\$ \allowbreak{},}). \\[2pt]
339 & 753 & \sichC{} & {(\cf Eq.~(\lbl{eq:b3-2.19}))} & {(\cf Part-II, Eq.~(\lbl{eq:b3-2.19}))} & Verweis & Schon in der Arbeitsliste. Gleiche Nummer (3.2.19) existiert auch in diesem Kapitel. \\[2pt]
340 & 782 & \sichA{} & $L_{p,o}$ & $L_{p,0}$ & Tippfehler & Buchstabe o statt Ziffer 0 (vgl. Eq. (3.2.23)). \\[2pt]
340 & 802 & \sichC{} & {$s(A)$ is strictly dominant provided that $\sigma(A) \neq \emptyset$} & {$s(L_{p})$ is strictly dominant provided that $\sigma(L_{p}) \neq \emptyset$} & Notation & Im Beispiel heißt der Generator $L_p$, A ist nicht eingeführt; ebenso Z. 803; bitte prüfen. \\[2pt]
340 & 814 & \sichC{} & $R(\lambda,A) : s(A) < \lambda$ & $R(\lambda,A) \colon s(A) < \lambda$ & Notation & Sonst steht in Mengenklammern durchgehend \code{\textbackslash{}colon} (andere Abstände); ebenso Z. 1163 ($\{\xi\in\C :\vert\xi\vert=1\}$). \\[2pt]
342 & 967 & \sichB{} & {$(b)\Rightarrow (a)$.} & {(b) $\Rightarrow$ (a).} & Notation & Im Mathemodus werden a, b kursiv gesetzt, anders als in Z. 575/593 und den Aufzählungslabels; ebenso Z. 974, 2079, 2089, 2284, 2291. \\[2pt]
344 & 1033 & \sichC{} & \code{\textbackslash{}begin\{align*\}\allowbreak{}\textbackslash{}label\{eq:c3-D0\}} & \code{\textbackslash{}begin\{align*\}} & LaTeX & \code{\textbackslash{}label} in einer unnummerierten Umgebung, nirgends referenziert; ebenso Z. 1056 (\code{\textbackslash{}begin\{equation*\}\allowbreak{}\textbackslash{}label\{eq:c3-r\_lambda\}}). \\[2pt]
345 & 1115 & \sichB{} & {Theorem~\lbl{thm:b2-3.5}} & {Part-II, Proposition~\lbl{prop:b3-3.5}} & Verweis & Schon in der Arbeitsliste. Gemeint ist laut „latter proposition“ und Beweis Z. 1140 Part-II, Proposition 3.3.5. \\[2pt]
345 & 1142 & \sichC{} & {by assertion (i) there exists} & {by assertion (ii) there exists} & Verweis & Schon in der Arbeitsliste. Strikte Positivität des Funktionals folgt aus (ii). \\[2pt]
346 & 1166 & \sichB{} & $\Gamma\times\Gamma\setminus(1,1)$ & $\Gamma\times\Gamma\setminus\{(1,1)\}$ & Notation & Mengendifferenz mit einem Punkt braucht Mengenklammern. \\[2pt]
346 & 1215 & \sichC{} & {Eq.~(\lbl{eq:a1-3.1})} & {Part-I, Eq.~(\lbl{eq:a1-3.1})} & Verweis & Gleichung aus Part I ohne Teilangabe, die Nummer (1.3.1) gibt es auch in Part III; ebenso Z. 1642 (eq:a2-1.9), 1717 (eq:a1-3.1), 2081 (eq:a1-3.2); vgl. Arbeitsliste zu Z. 753. \\[2pt]
348 & 1309 & \sichB{} & $C(K) \cong E_{h}$ & $C(K) \cong E_{\vert h\vert}$ & Tippfehler & h ist nicht positiv; in Z. 1306 wird $E_{\vert h\vert}$ mit C(K) identifiziert. \\[2pt]
348 & 1309 & \sichC{} & {We define $j \colon C(\Gamma) \to E$} & {We define $i \colon C(\Gamma) \to E$} & Notation & Schon in der Arbeitsliste. In der Aussage (Z. 1283) heißen die Abbildungen i und j, im Beweis j und $\tilde{j}$; ebenso Z. 1310, 1313, 1315, 1317, 1318. \\[2pt]
348 & 1313 & \sichC{} & $\int_{0}^{2\pi} e_{n}(\eu^{\im t})\dt  =2\pi\cdot\delta_{n0}$ & $\frac{1}{2\pi}\int_{0}^{2\pi} e_{n}(\eu^{\im t})\dt = \delta_{n0}$ & Tippfehler & Bitte prüfen: Wegen $\ensuremath{\langle}\vert h\vert,\ensuremath{\varphi}\ensuremath{\rangle} = 1$ (Z. 1305) müsste $\ensuremath{\langle}h^{[0]},\ensuremath{\varphi}\ensuremath{\rangle} = 1$ sein, nicht $2\ensuremath{\pi}$; ebenso Z. 1315. \\[2pt]
348 & 1342 & \sichB{} & $\alpha \notin P\sigma(A)$ & $\alpha \notin P\sigma(R)$ & Tippfehler & Im Lemma kommt nur der Operator R vor (vgl. Z. 1336). \\[2pt]
349 & 1356 & \sichC{} & $\mathcal{F}\text{-}\lim\sup$ & $\mathcal{F}\text{-}\limsup$ & Notation & \code{\textbackslash{}lim\allowbreak{}\textbackslash{}sup} setzt zwei getrennte Operatoren; üblich ist \code{\textbackslash{}limsup}. \\[2pt]
350 & 1436 & \sichB{} & $\vert\hat{f}\vert = \gamma R(\gamma+\im\beta,A)_{\mathcal{U}}\vert\hat{f}\vert$ & $\vert\hat{f}\vert = \vert\gamma R(\gamma+\im\beta,A)_{\mathcal{U}}\hat{f}\vert$ & Klammern & Bitte prüfen: Die Betragsstriche müssen um den ganzen Ausdruck stehen (Eigenvektor \code{\textbackslash{}hat\{f\} \allowbreak{}\ensuremath{\in} \allowbreak{}M}), erst dann folgt die Abschätzung. \\[2pt]
350 & 1443 & \sichC{} & $\hat{\phi}(\vert\hat{f}\vert) = P_{\mathcal{U}}\vert\hat{f}\vert$ & $\hat{\phi}(\vert\hat{f}\vert)\hat{u} = P_{\mathcal{U}}\vert\hat{f}\vert$ & Tippfehler & Bitte prüfen: Links steht ein Skalar, rechts ein Vektor (\code{P\_U \allowbreak{}= \allowbreak{}\allowbreak{}\textbackslash{}hat\{\ensuremath{\varphi}\}\ensuremath{\otimes}\allowbreak{}\textbackslash{}hat\{u\}}). \\[2pt]
350 & 1449 & \sichB{} & $\tilde{\phi}(\vert\hat{f}\vert) - \vert\hat{f}\vert \in I$ & $\hat{\phi}(\vert\hat{f}\vert)\hat{u} - \vert\hat{f}\vert \in I$ & Tippfehler & $\tilde{\ensuremath{\varphi}}$ ist nicht definiert (gemeint \code{\textbackslash{}hat\{\ensuremath{\varphi}\}}); zudem fehlt \code{\textbackslash{}hat\{u\}} (vgl. Z. 1443 und 1455); bitte prüfen. \\[2pt]
350 & 1453 & \sichB{} & $E_{\mathcal{U}/{I}}$ & $E_{\mathcal{U}}/I$ & Notation & Der Quotient wird im Index gesetzt; vgl. $E_1/I$ in Z. 690. \\[2pt]
350 & 1467 & \sichA{} & $J_{1},J_{2}, .. J_{m}$ & $J_{1},J_{2},\ldots,J_{m}$ & Tippfehler & Zwei statt drei Punkte; ebenso Z. 1473 (\code{J\_\{2\}\allowbreak{}\textbackslash{}oplus \allowbreak{}.. \allowbreak{}\allowbreak{}\textbackslash{}oplus \allowbreak{}J\_\{m\} \allowbreak{}.\$} $\ensuremath{\to}$ \code{J\_\{2\}\allowbreak{}\textbackslash{}oplus \allowbreak{}\allowbreak{}\textbackslash{}cdots \allowbreak{}\allowbreak{}\textbackslash{}oplus \allowbreak{}J\_\{m\}\$.}). \\[2pt]
351 & 1504 & \sichB{} & {To show that $T_{k}$ is irreducible} & {To show that $\TT_{k}$ is irreducible} & Notation & Die eingeschränkte Halbgruppe heißt $\TT_k$ (Z. 1468, 1472). \\[2pt]
351 & 1533 & \sichB{} & {Moreover, $\sigma_{b}$ contains} & {Moreover, $\sigma_{b}(A)$ contains} & Notation & Argument fehlt. \\[2pt]
352 & 1545 & \sichC{} & $\gamma_{1} \coloneq  \pi/v_{1}$, $\gamma_{2} \coloneq  \pi/v_{2}$ & $\gamma_{1} \coloneq \pi v_{1}$, $\gamma_{2} \coloneq \pi v_{2}$ & Tippfehler & Bitte prüfen: Wegen $v_1 < v_2$ ist so $\ensuremath{\gamma}_1 > \ensuremath{\gamma}_2$ und die Bedingung $n\ensuremath{\gamma}_1 \ensuremath{\leq} \vert \ensuremath{\gamma}\vert \ensuremath{\leq} n\ensuremath{\gamma}_2$ in Z. 1543 für $n \ensuremath{\geq} 1$ leer. \\[2pt]
352 & 1552 & \sichC{} & \code{\textbackslash{}textit\{\allowbreak{}\textbackslash{}small \allowbreak{}\$ \allowbreak{}(T(t)f)(x,v) \allowbreak{}= \allowbreak{}\allowbreak{}\textbackslash{}begin\{cases\} \allowbreak{}\dots{} \allowbreak{}\$\}} & $(T(t)f)(x,v) = \begin{cases} \dots  \end{cases}$ & LaTeX & Die Formel wird als kleine Inline-Formel in kursivem Text gesetzt, „if“ ist kursiv (anders als Eq. (3.3.3)); außerdem {k \textit{ even}} neben $k \text{ odd}$; ebenso Z. 1649. \\[2pt]
352 & 1570 & \sichA{} & $R(.,A)$ & $R(\cdot,A)$ & Notation & Sonst \code{\textbackslash{}cdot}, z. B. Z. 1576 und 1577. \\[2pt]
354 & 1664 & \sichA{} & $\dim(\ker((\eu^{\tau\cdot s(B)} - S(\tau))) > 1$ & $\dim(\ker(\eu^{\tau\cdot s(B)} - S(\tau))) > 1$ & Klammern & Eine öffnende Klammer zu viel. \\[2pt]
354 & 1675 & \sichA{} & $\dx'dv'$ & $\dx'\diff{v'}$ & Notation & Das zweite Differential hat ein kursives d. \\[2pt]
354 & 1678 & \sichB{} & {(\cf Eq.~(\lbl{eq:c3-3.5}) and Eq.~(\lbl{eq:c3-3.7}))} & {(\cf Eq.~(\lbl{eq:c3-3.5}) and Eq.~(\lbl{eq:c3-3.6}))} & Verweis & Bitte prüfen: $\ensuremath{\kappa}$ wird in Eq. (3.3.6) eingeführt, Eq. (3.3.7) ist das Irreduzibilitätskriterium. \\[2pt]
354 & 1679 & \sichA{} & {$K$ is A-compact} & {$K$ is $A$-compact} & Mathemodus & A steht ohne Mathemodus (vgl. Z. 1608, 1680). \\[2pt]
355 & 1717 & \sichB{} & $(\eu^{rt} - T(r))R(r,A) \geq 0$ & $(\eu^{rt} - T(t))R(r,A) \geq 0$ & Tippfehler & Bitte prüfen: Die folgende Abschätzung betrifft T(t). \\[2pt]
355 & 1749 & \sichB{} & $A^*(q_{\psi}'(x)) = q_{\psi}'(x)$ & $A^*(q_{\psi}'(x)) = \lambda q_{\psi}'(x)$ & Tippfehler & Bitte prüfen: Der Eigenwert $\ensuremath{\lambda}$ fehlt (vgl. vorangehende Gleichungen). \\[2pt]
357 & 1834 & \sichA{} & $P\sigma(A)\cap i\R$ & $P\sigma(A)\cap \im\R$ & Notation & Kursives i statt \code{\textbackslash{}im}; ebenso Z. 1869 ($iF\cap E_{\R }$), 1902 (zweimal $i\R$), 2107 ($\mu+i\R$). \\[2pt]
357 & 1836 & \sichA{} & \code{\textbackslash{}item \allowbreak{}E \allowbreak{}is \allowbreak{}weakly \allowbreak{}sequentially \allowbreak{}complete} & \code{\textbackslash{}item \allowbreak{}\$E\$ \allowbreak{}is \allowbreak{}weakly \allowbreak{}sequentially \allowbreak{}complete} & Mathemodus & E ohne Mathemodus. \\[2pt]
357 & 1848 & \sichA{} & $\im\alpha\Z \in P\sigma(A_{\vert F})$ & $\im\alpha\Z \subseteq P\sigma(A_{\vert F})$ & Tippfehler & Inklusion von Mengen. \\[2pt]
357 & 1853 & \sichC{} & {$T(t)'\phi = \phi$ and $\phi \gg 0 ,  (t \geq 0)$} & {$T(t)'\phi = \phi$ $(t \geq 0)$ and $\phi \gg 0$} & Satzzeichen & Komma und Bedingung ($t \ensuremath{\geq} 0$) stehen falsch innerhalb der Formel. \\[2pt]
358 & 1870 & \sichA{} & $F\cap\mathbb{E}_{\R }$ & $F\cap E_{\R }$ & Tippfehler & $\mathbb{E}$ statt E (vgl. Z. 1869). \\[2pt]
358 & 1897 & \sichC{} & $\int_{0}^{2\pi} f(z \cdot \eu^{\im s}) \ds$ & $\int_{0}^{2\pi} f(\eu^{\im s}) \ds$ & Tippfehler & Bitte prüfen: In der Definition von $g_t(x)$ ist z frei. \\[2pt]
358 & 1928 & \sichA{} & $r \in R(\sigma(T(t)))$ & $r \in R\sigma(T(t))$ & Notation & Residualspektrum, sonst als $R\sigma$ geschrieben (Z. 1930). \\[2pt]
360 & 2008 & \sichB{} & $\text{ for every } \ f \in D(A) ,$ & $\text{ for every } \ f \in D(A) .$ & Satzzeichen & Nach der Formel beginnt ein neuer Satz („For real $\ensuremath{\mu}$ \dots{}“). \\[2pt]
362 & 2204 & \sichB{} & $r(T(-t)_{\vert J_{\mu}}) < \eu^{\mu t}$ & $r(T(-t)_{\vert J_{\mu}}) < \eu^{-\mu t}$ & Tippfehler & Bitte prüfen: Nur so folgt Z. 2213 ($\vert \ensuremath{\lambda}\vert > e^{\ensuremath{\mu}t}$ nach Inversion). \\[2pt]
363 & 2260 & \sichB{} & \code{\textbackslash{}begin\{enumerate\} \allowbreak{}[(a)]} & \code{\textbackslash{}begin\{enumerate\}[\allowbreak{}\textbackslash{}upshape \allowbreak{}(a)]} & LaTeX & In der Proposition werden die Labels sonst kursiv gesetzt; alle anderen Aufzählungen verwenden \code{\textbackslash{}upshape}. \\[2pt]
363 & 2286 & \sichB{} & $(\sup\{\|T(t)\| \colon 0 \leq t \leq t_0\})^{1/p}$ & $(\sup\{\|T(t)\| \colon 0 \leq t \leq t_0\})^{p}$ & Tippfehler & Bitte prüfen: Z. 2289 braucht $\|T(t)\|^p \ensuremath{\leq} M$ und Z. 2303/2304 $\|T(t)\| \ensuremath{\leq} M^{1/p}$. \\[2pt]
364 & 2356 & \sichA{} & {on $L(\R ^n)$ $(1 \leq p < \infty)$} & {on $L^p(\R ^n)$ $(1 \leq p < \infty)$} & Tippfehler & Exponent p fehlt. \\[2pt]
365 & 2369 & \sichA{} & {n$\times$n-matrix} & {$n\times n$-matrix} & Mathemodus & n steht aufrecht außerhalb des Mathemodus. \\[2pt]
365 & 2411 & \sichC{} & {Part-III, Section~\lbl{sec:c4-3}} & {Section~\lbl{sec:c4-3}} & Verweis & Innerhalb von Part III sonst ohne Teilangabe zitiert (vgl. Arbeitsliste Kap. 4, S. 372). \\
\end{longtable}
\subsection*{Teil III, Kap. 4 – Asymptotics of Positive Semigroups on Banach Lattices (\code{part-c/chap-c4.tex})}
\begin{longtable}{@{}rrc>{\RaggedRight\arraybackslash}p{5.8cm}>{\RaggedRight\arraybackslash}p{5.8cm}>{\RaggedRight\arraybackslash}p{1.9cm}>{\RaggedRight\arraybackslash}p{7.3cm}@{}}
\toprule
\textbf{S.} & \textbf{Z.} & \textbf{} & \textbf{Ist} & \textbf{Soll} & \textbf{Art} & \textbf{Kommentar} \\
\midrule
\endfirsthead
\toprule
\textbf{S.} & \textbf{Z.} & \textbf{} & \textbf{Ist} & \textbf{Soll} & \textbf{Art} & \textbf{Kommentar} \\
\midrule
\endhead
\bottomrule
\endlastfoot
369 & 27 & \sichA{} & {Part-III, Chapter~\lbl{chap:b3}} & {Chapter~\lbl{chap:c3}} & Verweis & Label zeigt auf Part-II, Kap. 3 (Hyperlink falsch, gedruckte Nummer zufällig gleich); gemeint ist Part-III, Kap. 3. \\[2pt]
369 & 55 & \sichB{} & {Remark~\lbl{rem:b4-1.5}.)} & {Remark~\lbl{rem:b4-1.5})} & Satzzeichen & Schon in der Arbeitsliste. Überflüssiger Punkt vor der Klammer. \\[2pt]
369 & 100 & \sichC{} & $L^{\,p}(\mu)$ & $L^{p}(\mu)$ & Abstand & Zusätzlicher Abstand im Exponenten; sonst überall $L^{p}$. \\[2pt]
371 & 154 & \sichB{} & $k \cdot \log [r(R(k,A_{1/u}))]^{-1}$ & $k \cdot \log [r(R(k,A_{1/u}))]$ & Tippfehler & Wegen $r(R(k,A_u)) = (k\ensuremath{-}s(u))^{-1}$ (Z. 141) ist $[k\ensuremath{-}s(1/u)]^{-1} = r(R(k,A_{1/u}))$; Exponent $\ensuremath{-}1$ überzählig – bitte prüfen. \\[2pt]
371 & 176 & \sichB{} & {Part-III, Theorems~\lbl{thm:c3-1.2}} & {Theorem~\lbl{thm:c3-1.2}} & Verweis & Schon in der Arbeitsliste. Plural bei nur einem Satz. \\[2pt]
373 & 243 & \sichB{} & $0 \notin \text{P}\sigma(A)$ & $0 \notin P\sigma(A)$ & Notation & \code{\textbackslash{}text\{P\}} setzt P im (im Theorem kursiven) Textfont; sonst überall mathematisches $P\sigma$ bzw. $P{\sigma}$. \\[2pt]
373 & 252 & \sichC{} & {\lbl{thm:c4-1.5-1} In case} & {\lbl{thm:c4-1.5-1}: In case} & Satzzeichen & Doppelpunkt nach der Implikation „$(b) \ensuremath{\Rightarrow} (a)$“ fehlt (vgl. Z. 274). \\[2pt]
373 & 277 & \sichC{} & $R{\sigma}(A^*) = P{\sigma}(A^{**})$ & $R{\sigma}(A') = P{\sigma}(A'')$ & Notation & Adjungierte sonst mit Strich (A', T(t)', Z. 234, 689, 969, 1251); ebenso Z. 278 $(T(t)^*)$, Z. 279 $A^*$. \\[2pt]
374 & 309 & \sichB{} & $E_{+} \subseteq \Image{A(D(A)_{-})}$ & $E_{+} \subseteq A(D(A)_{-})$ & Notation & \code{\textbackslash{}Image} setzt „im(\dots{})“ um eine Menge, die bereits das Bild ist – bitte prüfen; ebenso Z. 316. \\[2pt]
374 & 341 & \sichB{} & $(\image{A})_{+}$ & $(\Image{A})_{+}$ & Notation & \code{\textbackslash{}image} ergibt „im A“, \code{\textbackslash{}Image} „im(A)“ – im selben Absatz gemischt; ebenso Z. 344, 355, 357 \code{(\allowbreak{}\textbackslash{}image)} gegenüber Z. 343, 348, 446 \code{(\allowbreak{}\textbackslash{}Image)}. \\[2pt]
374 & 356 & \sichC{} & $T(t)f \, \dt$ & $T(t)f \dt$ & Abstand & \code{\textbackslash{}dt/\allowbreak{}\textbackslash{}diff} bringen den Abstand schon mit, mit \code{\textbackslash{}}, doppelt (Kapitel sonst meist ohne); ebenso u. a. Z. 343, 394, 452, 684, 897, 921, 925, 1214, 1229, 1241, 1248, 1260, 1794, 1801, 1884. \\[2pt]
374 & 357 & \sichB{} & \code{\$\allowbreak{}\textbackslash{}sigma(E',E)\$-\$\allowbreak{}\textbackslash{}lim} & \code{\$\allowbreak{}\textbackslash{}sigma(E,E')\$-\$\allowbreak{}\textbackslash{}lim} & Tippfehler & Schwache Topologie auf E ist $\ensuremath{\sigma}(E,E')$ (vgl. Z. 382) – bitte prüfen. \\[2pt]
375 & 379 & \sichC{} & $\lim_{\lambda \downarrow 0}$ & $\lim_{\lambda \to 0+}$ & Notation & Sonst im Kapitel $\ensuremath{\lambda} \ensuremath{\to} 0+$ (Z. 362, 368, 389). \\[2pt]
376 & 435 & \sichB{} & $g_{0} \in \Image{A_{+}}$ & $g_{0} \in (\Image{A})_{+}$ & Notation & „$\operatorname{im}(A_+)$“ ergibt keinen Sinn; gemeint ist der positive Teil des Bildes (vgl. Z. 341, 452). \\[2pt]
376 & 435 & \sichB{} & $\lim_{t \to \infty} u(t) = -h$ & $\lim_{t \to \infty} u(t) = h$ & Tippfehler & Im Beweis (Z. 446) konvergiert u(t) gegen h mit $Ah = \ensuremath{-}f_0$ (vgl. Prop. 1.12: $\ensuremath{-}A^{-1}f_0$) – Vorzeichen bitte prüfen. \\[2pt]
379 & 545 & \sichC{} & {Example~\lbl{ex:a3-1.3} \&~\lbl{ex:a3-1.4}} & {Examples~\lbl{ex:a3-1.3} and~\lbl{ex:a3-1.4}} & Verweis & Schon in der Arbeitsliste. \\[2pt]
379 & 560 & \sichB{} & {Section~\lbl{sec:b4-2}} & {Part-II, Section~\lbl{sec:b4-2}} & Verweis & Schon in der Arbeitsliste. Ohne Part-Angabe liest sich „Section 4.2“ als der aktuelle Abschnitt. \\[2pt]
379 & 577 & \sichB{} & {Eq.~(\lbl{eq:b4-2.5})} & {Part-II, Eq.~(\lbl{eq:b4-2.5})} & Verweis & Gleichung aus Part-II trägt dieselbe Nummer (4.2.5) wie Eq. (4.2.5) dieses Kapitels (Z. 752); Part-Angabe fehlt. \\[2pt]
380 & 597 & \sichC{} & $\omega_{ess}(\TT)$ & $\omega_{\mathrm{ess}}(\TT)$ & Mathemodus & „ess“ kursiv wie ein Produkt gesetzt; ebenso Z. 645, 669, 1178 ($r_{ess}$). Buchweit so – nur global ändern. \\[2pt]
381 & 661 & \sichC{} & $Af  := -f' - \mu f$ & $Af  \coloneq -f' - \mu f$ & Notation & Sonst durchgehend \code{\textbackslash{}coloneq} (zentrierter Doppelpunkt); ebenso Z. 662, 1339, 1881. \\[2pt]
383 & 830 & \sichA{} & {assertion~\lbl{thm:c4-2.6-1} , then} & {assertion~\lbl{thm:c4-2.6-1}, then} & Satzzeichen & Leerzeichen vor Komma nach \code{\textbackslash{}ref}; ebenso Z. 911 ({\lbl{ex:c4-2.7-1} ,}). \\[2pt]
384 & 867 & \sichA{} & {$t > 0$ ,} & {$t > 0$,} & Satzzeichen & Leerzeichen vor Komma; ebenso Z. 871 ({$t > 0$\, ,}), 1017 ({$L^p(\mu)$ ,}), 1494 ({$h \colon \R \to \R_{+}$ ,}). \\[2pt]
384 & 875 & \sichB{} & $\int_{K} P_{s}(y,C)P_{t}(x,dy)$ & $\int_{X} P_{s}(y,C)P_{t}(x,\diff{y})$ & Tippfehler & Grundmenge heißt X (Z. 861), K ist nicht definiert; ebenso $x \in K$ $\ensuremath{\to}$ $x \in X$ in derselben Zeile; dazu kursives d – bitte prüfen. \\[2pt]
384 & 883 & \sichB{} & $\, d\mu(x)$ & $\, \diff{\mu(x)}$ & Notation & Kursives d; sonst \code{\textbackslash{}diff} (Z. 897, 921, 925); ebenso Z. 906 (zweimal). \\[2pt]
384 & 892 & \sichB{} & $= 1_{C}(x)$ & $= \1_{C}(x)$ & Notation & Indikatorfunktion sonst fett $\1$ (Z. 787, 917, 925). \\[2pt]
384 & 892 & \sichB{} & \code{\textbackslash{}mu\allowbreak{}\textbackslash{}text\{-a.e.\}.\$} & \code{\textbackslash{}mu\allowbreak{}\textbackslash{}text\{-a.e.\}\$} & Satzzeichen & Doppelter Punkt „a.e..“. \\[2pt]
385 & 917 & \sichB{} & $\quad\mu-\text{a.e.}.$ & $\quad\mu\text{-a.e.}$ & Mathemodus & Bindestrich im Mathemodus wird zum Minuszeichen; dazu doppelter Punkt. \\[2pt]
385 & 970 & \sichA{} & {Since Pe = e, hence} & {Since $Pe = e$, hence} & Mathemodus & Formel ohne $\dots $. \\[2pt]
386 & 1003 & \sichC{} & $T^n \wedge T^m > 0$ & $\inf\{T^n,T^m\} > 0$ & Notation & Infimum von Operatoren sonst als \code{\textbackslash{}inf\allowbreak{}\textbackslash{}\{\allowbreak{}\textbackslash{}cdot,\allowbreak{}\textbackslash{}cdot\allowbreak{}\textbackslash{}\}} (Z. 957, 961). \\[2pt]
386 & 1021 & \sichA{} & $\cap \im\R = {0}$ & $\cap \im\R = \{0\}$ & Klammern & Unmaskierte Mengenklammern: gedruckt wird „$= 0$“. \\[2pt]
388 & 1088 & \sichC{} & $T = T_{1} \oplus T_{2}$ & $\TT = \TT_{1} \oplus \TT_{2}$ & Notation & T ohne Argument nicht definiert; vorher (Z. 1067) $\TT_1$, $\TT_2$. \\[2pt]
389 & 1151 & \sichB{} & $Mh \coloneq h\cdot f$ & $Mf \coloneq h\cdot f$ & Tippfehler & M wirkt auf $f \ensuremath{\in} L^p(\ensuremath{\nu})$, h ist die feste Dichte – bitte prüfen. \\[2pt]
389 & 1159 & \sichB{} & $P = h \otimes l$ & $P = h \otimes \1$ & Tippfehler & Nach der Normierung $T(t)\ensuremath{\mathbf{1}} = \ensuremath{\mathbf{1}}$ ist $\ensuremath{\mathbf{1}}$ der Fixvektor; „l“ vermutlich verlesene 1 – bitte prüfen. \\[2pt]
389 & 1172 & \sichB{} & $\|R(1,A)^{n} - P\|^{2/p}\|\cdots\|^{1-2/p}$ & $\|R(1,A)^{n} - P\|^{p/2}\|\cdots\|^{1-p/2}$ & Tippfehler & Für $p < 2$ wäre $1\ensuremath{-}2/p$ negativ; Riesz-Konvexität (vgl. Z. 1155: Exponent p/q) liefert p/2 und $1\ensuremath{-}p/2$ – bitte prüfen; ebenso Z. 1173. \\[2pt]
389 & 1176 & \sichB{} & $R(1,A)(1-P)$ & $R(1,A)(\Id-P)$ & Notation & Identität unmittelbar vorher (Z. 1165) als \code{\textbackslash{}Id}. \\[2pt]
390 & 1202 & \sichC{} & $f(\frac{1}{2}) = 0$ & $f(\frac{1}{4}) = 0$ & Tippfehler & Intervall ist [1/4,1]; Halbgruppe (Z. 1223–1226) und Resolvente (Z. 1229) passen zur Randbedingung bei 1/4 – bitte prüfen. \\[2pt]
390 & 1240 & \sichB{} & $\text{\sfrac{1}{x}}$ & $\mfrac{1}{x}$ & Notation & \code{\textbackslash{}sfrac} ist das Textmakro, x erscheint aufrecht statt kursiv; ebenso Z. 1241. \\[2pt]
392 & 1298 & \sichB{} & $\Phi(u_t),  t \geq 0$ & $\Phi(u_t), \quad t \geq 0$ & Abstand & Leerzeichen im Mathemodus wirkungslos, „$t \ensuremath{\geq} 0$“ klebt am Komma; ebenso Z. 1315, 1794 ($t \geq 0, f_{0}$). \\[2pt]
393 & 1349 & \sichB{} & {Example~\lbl{subsec:a1-2.4}.\,(ii)} & {Example~\lbl{subsec:a1-2.4}\,(ii)} & Satzzeichen & Schon in der Arbeitsliste. Überflüssiger Punkt. \\[2pt]
393 & 1381 & \sichC{} & {Formula Eq.~(\lbl{eq:c4-3.2})} & {Eq.~(\lbl{eq:c4-3.2})} & Verweis & Doppelte Bezeichnung „Formula Eq. (\dots{})“; ebenso Z. 1456 („formula Eq.“), 1473 („equation Eq.“), 1533, 1672. \\[2pt]
393 & 1384 & \sichC{} & $\|\sum_{n=0}^{\infty} S_{\lambda}^n\|$ & $\Bigl\|\sum_{n=0}^{\infty} S_{\lambda}^n\Bigr\|$ & Klammern & Normstriche in der abgesetzten Formel nicht an die Summe angepasst. \\[2pt]
394 & 1433 & \sichC{} & {supp $g \subseteq [-1,-t]$} & $\supp g \subseteq [-1,-t]$ & Mathemodus & „supp“ im Fließtext; Makro \code{\textbackslash{}supp} vorhanden. \\[2pt]
394 & 1463 & \sichA{} & $\|w_{t}(s)\_{F}\|$ & $\|w_{t}(s)\|_{F}$ & LaTeX & $\_$ druckt einen Unterstrich; Normindex falsch platziert. \\[2pt]
394 & 1474 & \sichC{} & {(use (\lbl{eq:c4-re}))} & {(use Eq.~(\lbl{eq:c4-3.1}))} & Verweis & Schon in der Arbeitsliste. Positivität folgt aus (4.3.1) – bitte prüfen. \\[2pt]
394 & 1475 & \sichA{} & {Proposition~\lbl{prop:c2-4.1} )} & {Proposition~\lbl{prop:c2-4.1})} & Satzzeichen & Schon in der Arbeitsliste (\sichA{}-Punkt S. 370, „ebenso“-Fälle), im Quelltext aber noch nicht korrigiert; ebenso Z. 1502 ({\lbl{cor:c3-2.12} )}). \\[2pt]
395 & 1484 & \sichC{} & $\Phi(\epsilon_{\lambda} \cdot x)$ & $\Phi(\epsilon_{\lambda} \otimes x)$ & Notation & Z. 1360–1363 definiert $\ensuremath{\varepsilon}_\ensuremath{\lambda} \ensuremath{\otimes} x$; so auch Z. 1522 und 1539. \\[2pt]
396 & 1549 & \sichC{} & {If \, $\lambda \in \sigma(A)$} & {If\/ $\lambda \in \sigma(A)$} & Abstand & Wortabstand plus \code{\textbackslash{}}, sichtbar zu breit; sonst $If\/$; ebenso Z. 1552. \\[2pt]
396 & 1574 & \sichC{} & {$\Phi(D(A_{0}))$ $= F$} & $\Phi(D(A_{0})) = F$ & LaTeX & Formel in zwei Mathe-Stücke zerrissen, Abstand vor „$=$“ falsch; ebenso Z. 1693. \\[2pt]
396 & 1574 & \sichC{} & $x - \Phi_{\lambda}(x) = 0$ & $x - \Phi_{\lambda}(x) = \Phi R(\lambda,A_{0})g$ & Tippfehler & Rechte Seite hängt nicht von g ab („for every $g \ensuremath{\in} E$“); gemeint vermutlich die Gleichung aus Z. 1540 – bitte prüfen. \\[2pt]
396 & 1576 & \sichC{} & {the proof of~\lbl{lem:c4-3.5}} & {the proof of Lemma~\lbl{lem:c4-3.5}} & Verweis & Schon in der Arbeitsliste. \\[2pt]
396 & 1576 & \sichB{} & {equivalent to $\lambda \in \sigma(A)$} & {equivalent to $\lambda \in \rho(A)$} & Tippfehler & Eindeutige Lösbarkeit $\ensuremath{\Leftrightarrow} 1 \ensuremath{\in} \ensuremath{\varrho}(\ensuremath{\Phi}_\ensuremath{\lambda}) \ensuremath{\Leftrightarrow} \ensuremath{\lambda} \ensuremath{\in} \ensuremath{\varrho}(A)$ (Lemma 3.5) – bitte prüfen. \\[2pt]
396 & 1578 & \sichB{} & $\sigma(\Phi_{\lambda})\backslash\{0\}$ & $\sigma(\Phi_{\lambda})\setminus\{0\}$ & Notation & Sonst \code{\textbackslash{}setminus} (Z. 1078, 1088); \code{\textbackslash{}backslash} ist ein anderes Zeichen ohne Operatorabstand. \\[2pt]
396 & 1586 & \sichC{} & $A \coloneq \frac{\diff{}}{\dt}$ & $A \coloneq \frac{\diff{}}{\dx}$ & Notation & Derselbe Operator heißt in Z. 1336 d/dx – bitte vereinheitlichen. \\[2pt]
396 & 1611 & \sichA{} & $s(\Phi_{r)} \geq 1$ & $s(\Phi_{r}) \geq 1$ & Klammern & Schließende Klammer im Index. \\[2pt]
397 & 1633 & \sichB{} & $h: \mu \mapsto s(\Phi_{\mu})$ & $h \colon \mu \mapsto s(\Phi_{\mu})$ & Notation & Sonst \code{\textbackslash{}colon} (Z. 1494, 1744, 1815, 1901); „:“ als Relation mit falschem Abstand. \\[2pt]
398 & 1681 & \sichA{} & $E = (L^{1}[-1,0],F)$ & $E = L^{1}([-1,0],F)$ & Tippfehler & „$L^1$“ fehlt vor der Klammer (vgl. Z. 1303, 1722). \\[2pt]
398 & 1712 & \sichB{} & $\psi(e_{\lambda})$ & $\psi(\epsilon_{\lambda})$ & Notation & Die Funktion heißt sonst $\ensuremath{\varepsilon}_\ensuremath{\lambda}$ (Z. 1360, 1696). \\[2pt]
398 & 1732 & \sichC{} & {Corollary~\lbl{cor:c4-3.2}, \ Eq.} & {Corollary~\lbl{cor:c4-3.2}, Eq.} & Abstand & Doppelter Wortabstand nach dem Komma. \\[2pt]
398 & 1775 & \sichB{} & {of\/ $K(s)$) ,} & {of\/ $K(s)$),} & Satzzeichen & Schon in der Arbeitsliste. \\[2pt]
398 & 1776 & \sichC{} & {Hence \ $\langle x_{\lambda},x_{\mu}'\rangle \neq 0$ \ and} & {Hence $\langle x_{\lambda},x_{\mu}'\rangle \neq 0$ and} & Abstand & Doppelte Wortabstände um die Formel. \\[2pt]
399 & 1801 & \sichB{} & $\dx' \ ,$ & $\dx',$ & Satzzeichen & Sichtbarer Abstand vor dem Komma in der abgesetzten Formel. \\[2pt]
399 & 1824 & \sichB{} & $\Phi_{\mu}(x) =$ & $\Phi_{\mu}f(x) =$ & Tippfehler & f fehlt (vgl. rechte Seite $\ensuremath{\Phi}_\ensuremath{\lambda}$ f(x), Z. 1827). \\[2pt]
401 & 1899 & \sichC{} & {{bellmancooke:1963}\,)} & \code{\{bellmancooke:1963\})} & Abstand & Schon in der Arbeitsliste (\sichA{}, S. 401); nach der Korrektur bleibt mit \code{\textbackslash{}}, ein kleiner Abstand vor „)“. \\[2pt]
401 & 1909 & \sichB{} & $(i=1,...,n)$ & $(i=1,\ldots,n)$ & LaTeX & „...“ im Mathemodus statt \code{\textbackslash{}ldots} (Abstände falsch); ebenso Z. 1913, 1917; Z. 2070 korrekt. \\[2pt]
402 & 2067 & \sichB{} & $\frac{\partial u_i}{\partial t}(t,a)$ & $\frac{\partial u}{\partial t}(t,a,x)$ & Tippfehler & System scheint aus Bsp. 3.12 (Z. 1909–1917) kopiert: n und $\ensuremath{\beta}_ij$ sind in Bsp. 3.13 nicht definiert, die Interpretation (Z. 2097) spricht von u(t,a,x), $\ensuremath{\beta}$ und $\ensuremath{\Omega}$. Ganze align-Umgebung Z. 2067–2079 bitte am Original prüfen. \\
\end{longtable}
\end{document}
